% MT059 / answer-en — canonical editable content.

\begin{enumerate}[start=1]

\item Both can yield linear decision boundaries in the given feature space when Gaussian NB shares variances across classes.

\item Both can be learned using maximum likelihood, but LR maximizes conditional label likelihood, whereas NB fits the joint model through class priors and class-conditional likelihoods.

\item NB is generative; LR directly models the class posterior and is discriminative.

\item Both have multiclass extensions and can incorporate regularization or parameter priors.

\item For binary classification with \(d\) features, LR has \(d+1\) coefficients, including the intercept. Under the key's \textbf{single shared scalar variance} convention, NB has \(2d+1\) class-conditional parameters: two mean vectors and one variance. A learned binary class prior adds one more free parameter, making \(2d+2\) in total.

\end{enumerate}


If “shared variance” means one shared variance \textbf{per feature}, NB has \(3d+1\) free parameters including the learned prior. State which model and counting convention you use.
